\documentclass[11pt,a4paper]{article}

\usepackage[T1]{fontenc}
\usepackage[utf8]{inputenc}
\usepackage[french]{babel}
\usepackage{lmodern}
\usepackage{microtype}
\usepackage[a4paper,margin=2.4cm]{geometry}
\usepackage{enumitem}
\usepackage[autostyle]{csquotes}
\usepackage{xurl}
\usepackage[hidelinks]{hyperref}
\usepackage{amsmath}
\usepackage[backend=biber,style=authoryear,maxcitenames=2,maxbibnames=12,pagetracker=false]{biblatex}
\addbibresource{../bibliographie/references.bib}

\setlength{\parindent}{0pt}
\setlength{\parskip}{0.6em}
\setlist[itemize]{leftmargin=1.5em,itemsep=0.25em,topsep=0.3em}
\setlist[enumerate]{leftmargin=1.8em,itemsep=0.35em,topsep=0.3em}

\newcommand{\projectversion}{14}
\newcommand{\projectdate}{6 septembre 2026}

\title{L'émergence de l'intéressant\\
\large Naissance des idées, compréhension et singularité des formes}
\author{François Pachet}
\date{Projet de thèse -- version courte, V\projectversion\\\small \projectdate}

\begin{document}

\maketitle

\begin{abstract}
Ce projet étudie l'intéressant comme objet philosophique. Il part d'une
expérience ordinaire : une idée surgit, un passage retient l'attention et nous
pousse à le relire ou à le réécouter pour comprendre ce qui nous a arrêtés.
L'intéressant se distingue du beau, de ce qui intrigue ou de ce qui surprend
par un trait précis : il déclenche un processus de construction mentale. Le
sujet commence alors à
distinguer, questionner, relier, modéliser, essayer ou produire autrement.
L'intéressant désigne le pouvoir de déclenchement qui naît dans la relation
entre un objet, un sujet, son histoire et une situation. Le mot
\emph{intérescence} nommera le processus ainsi ouvert, de la rupture avec
l'indifférence jusqu'à l'abandon, à une nouvelle construction ou à la
transmission. La thèse partira de Christian Garve, qui a formulé dès 1779 une
théorie de l'attention, de la lacune et de l'activité des représentations. Elle
étudiera ensuite les propriétés de l'intéressant dans trois domaines : la
musique, l'intelligence artificielle et la littérature. La créativité sera une
question commune à ces domaines ; la philosophie fournira l'approche générale.
La thèse examinera enfin une hypothèse nouvelle : l'intéressant peut contribuer à faire
naître un objet pour un sujet ; la réification commence lorsque l'histoire de
cette relation est oubliée.
\end{abstract}

\section{Origine et question}

Ce travail est issu d'une pratique constructiviste et empirique de la recherche
en intelligence artificielle. J'ai passé quarante ans à construire des systèmes
pour générer, transformer ou accompagner des musiques et des textes. Le projet
doctoral part d'une question personnelle : \emph{comprendre ce que je fais}
lorsque je construis un système, reconnais une solution, élimine une production
correcte mais sans intérêt, ou m'arrête devant une forme qui semble s'imposer.

La question directrice est la suivante :

\begin{quote}
Dans quelles conditions un objet déclenche-t-il chez un sujet un processus de
construction ? Comment ce processus se poursuit-il, s'épuise-t-il ou se
transforme-t-il ?
\end{quote}

Une chanson, une démonstration, une phrase ou une idée peut provoquer des
comparaisons, donner envie de la relire ou de la réécouter et susciter des
variations. Cet effet ne se réduit pas à la nouveauté : l'aléatoire surprend
sans nécessairement intéresser. Il ne se réduit pas non plus à la beauté, à la
popularité ou à une préférence donnée une fois pour toutes. Le travail devra
donc situer l'intéressant parmi ces notions voisines, en donner une définition
assez précise pour construire des contre-exemples, puis éprouver cette
définition dans des pratiques effectives.

\section{Un problème philosophique dispersé}

Le mot \emph{intérêt} vient du latin \emph{interest} : « il importe » et « il y
a une différence ». Le français conserve plusieurs sens : porter intérêt à une
question, avoir intérêt à agir, défendre ses intérêts ou prêter à intérêt
\parencite{academie2026interet,cnrtl2026interet}. Dans tous ces cas, quelque
chose compte pour un sujet, mais ce « compter » peut désigner une attention, un
enjeu pratique, un avantage ou une valeur monétaire. La thèse porte sur le
premier sens tout en gardant à l'esprit cette histoire commune.

La philosophie a étudié plusieurs dimensions de l'intéressant sous d'autres
noms : jugement, invention, mémoire, attente, curiosité, surprise, ennui ou
changement d'aspect. Chacun de ces concepts apporte une distinction utile. Le
vrai évalue une proposition ; le beau relève d'un jugement esthétique ; la
surprise marque un écart à l'attente ; la curiosité vise souvent la réponse à
une question déjà formulée. Aucun n'explique seul pourquoi une forme met un
sujet au travail et transforme ses manières de percevoir, de questionner ou
d'agir.

L'intéressant pourrait nommer cette fonction. On peut aimer un morceau sans le
trouver intéressant, être surpris sans rien construire, reconnaître une
proposition vraie sans vouloir la poursuivre, ou être mobilisé par un objet
pénible et moralement condamnable. Il ne s'agit donc pas d'inventer un domaine
entièrement neuf, mais de réunir un problème resté dispersé.

\section{Garve : un antécédent direct}

Dans \emph{Einige Gedanken über das Interessirende}, Christian Garve décrit en
1779 la manière dont un objet fixe l'attention sans effort volontaire
\parencite[253--274, 313--317]{garve1779interessirende}. L'attraction dépend du
rapport entre l'objet et l'état antérieur du sujet. Elle se maintient tant que
la suite conserve une part d'obscurité.

Garve va plus loin. Une explication ne devient intéressante que si le sujet a
d'abord éprouvé, même confusément, la difficulté qu'elle résout. L'idée nouvelle
éclaire des pensées déjà présentes mais incomplètes, répare un manque rencontré
dans une opération ou se construit avec des matériaux antérieurs. Lorsque
l'issue devient certaine, l'activité diminue. On trouve donc déjà chez Garve la
lacune aperçue, la dépendance à l'histoire du sujet, le désir de poursuivre, la
construction et l'épuisement local de l'intérêt.

La définition proposée ici ne sera pas présentée comme une découverte sans
précédent. Son apport résidera dans une description plus systématique du
processus, dans la distinction entre lacune, énigme, problème et contrainte,
dans le passage de la réception à la création et dans l'étude de la
récursivité. La reconstruction de \textcite{nannini2018garve} et le lexique de
\textcite{nannini2018interesting} guideront ce retour au texte allemand.

Ce travail donnera aussi lieu à un projet éditorial mené avec Christian Berner
pour les éditions Vrin : une traduction française intégrale de l'essai de
Garve, accompagnée d'une introduction et de notes. La traduction fait partie de
la recherche, car les mots de l'intérêt, de l'attention, de la lacune et de
l'activité mentale engagent directement l'interprétation du texte.

\subsection{Pourquoi Garve n'a-t-il pas fait école ?}

La réponse demandera une enquête historique. Une première cause possible est le
déclassement de la \emph{Popularphilosophie}. Garve défendait une philosophie
empirique, publique et utile, progressivement rejetée au profit de systèmes qui
ont acquis une plus grande autorité \parencite{thompson2018garve,
bachmannmedick2008anziehungskraft}. Une deuxième cause est le déplacement de la
question vers l'esthétique : Kant organise le jugement de goût autour du
désintéressement, puis Schlegel fait de l'intéressant un caractère de la poésie
moderne opposé à la beauté antique \parencite{kant1993juger,
schlegel2001study}. Une troisième cause est la dispersion du problème entre
psychologie, sciences cognitives, esthétique, épistémologie et sciences de
l'éducation.

Quelques auteurs ont pourtant étudié directement l'intéressant. Stace, Kolnai,
Ngai, Epstein et Grimm ont analysé sa valeur esthétique, discursive ou
épistémique \parencite{stace1944interestingness,kolnai1964concept,
ngai2008merely,epstein2009interesting,grimm2011interesting}. Plus récemment,
\textcite{vazard2026inquiry} a distingué la curiosité, dirigée vers une question,
de l'intérêt, dirigé vers un objet par une attention large et exploratoire.
L'intérêt peut ainsi précéder la question qu'il contribuera à faire apparaître.
Ces travaux montrent une histoire discontinue, sans théorie commune comparable
à celles du vrai ou du beau.

\subsection{Le risque de remplacer la question par une mesure}

La recherche scientifique rend un problème traitable en le définissant, en le
décomposant et en choisissant des mesures. Cette opération est indispensable,
mais elle peut déplacer la question sans que ce déplacement soit aperçu.
« Produire quelque chose d'intéressant » devient alors produire quelque chose
de correct, de cohérent, de surprenant, de probable dans un corpus ou de
conforme à des contraintes. Ces propriétés sont mesurables ; aucune ne suffit à
décrire ce qu'un objet fait construire à un sujet.

La thèse distinguera donc trois éléments : la cible initiale, l'indicateur choisi
pour l'approcher et la procédure qui calcule ou optimise cet indicateur. Un bon
score ne prouve pas que la cible a été atteinte. Inversement, un cas qui résiste
à la mesure ne doit pas être éliminé pour cette seule raison. L'avertissement de
\textcite{russell1995rationality} contre la mathématisation prématurée guidera
ce travail : la précision gagnée sur un substitut peut faire perdre le problème
que l'on voulait comprendre.

La parcimonie tonale offre un exemple positif
\parencite{pachet2026tonalparsimony}. L'analyse cherche d'abord des transitions
harmoniques continues, puis un petit nombre de centres tonals. Certains passages
résistent pourtant à cette compression : une analyse plus courte peut effacer
une excursion ou deux fonctions volontairement distinctes. Ces cas ne sont pas
écartés comme des anomalies ; ils montrent où le modèle doit être enrichi. Une
formalisation utile doit ainsi conserver la trace de ce qu'elle ne saisit pas.

\section{Définition et concepts principaux}

Le noyau de la thèse peut être formulé ainsi :

\begin{quote}
\textbf{Est intéressant, pour un sujet, ce qui déclenche chez lui un processus
de construction.}
\end{quote}

Construire ne signifie pas nécessairement fabriquer une chose. Le sujet
construit lorsqu'il modifie sa manière de percevoir, de questionner, de
comprendre ou d'agir : il fait une distinction, formule une question, relie des
éléments, élabore une hypothèse, essaie une variation ou invente une opération.
Une \emph{prise} est précisément l'un de ces résultats réutilisables : une
distinction, une question, une comparaison, un modèle ou une opération qui
permet de poursuivre. La simple exécution d'une procédure déjà maîtrisée ne
suffit pas.

La définition porte sur le déclenchement, pas sur la réussite. Le processus peut
être bref, échouer, s'égarer, épuiser son objet ou produire une capacité
nouvelle. Elle n'attribue pas non plus l'intéressant à l'objet seul. Une même
forme peut intéresser un sujet et laisser un autre indifférent ; elle peut aussi
cesser d'intéresser après apprentissage. Cette relation peut être notée
\[
  I(F,S\mid H,t),
\]
où \(F\) désigne la forme rencontrée, \(S\) le sujet, \(H\) son horizon social,
culturel et technique, et \(t\) le moment. La notation rappelle les variables
dont dépend le déclenchement ; elle ne remplace pas l'expérience par une mesure.

\subsection{Quelle universalité ?}

La dépendance au sujet ne condamne pas l'intéressant à une préférence privée
impossible à discuter. L'universalité recherchée n'est pas une propriété de
l'objet : aucun morceau, texte ou problème ne doit intéresser tout le monde. Elle
ne se trouve pas non plus dans le sujet pris isolément, puisque son histoire, ses
compétences et la situation de la rencontre comptent.

Kant fournit ici un point de comparaison précis. Dans les paragraphes 6 à 8 de
la \emph{Critique de la faculté de juger}, le jugement de goût ne repose pas sur
un concept de l'objet, mais prétend pourtant à une « universalité subjective »
\parencite[§§ 6--8]{kant1993juger}. Le jugement sur l'intéressant ne reprend pas
cette prétention : dire qu'une chose est intéressante ne demande pas que chacun
éprouve le même intérêt pour elle.

Ce qui peut être universel est la forme de la relation : lorsqu'un objet devient
intéressant, il se détache pour un sujet situé et déclenche une construction.
L'universel est donc ici celui d'une relation, non celui d'un objet ou d'un goût.
Cette hypothèse permet la discussion : on peut donner des raisons de poursuivre
l'examen, décrire les constructions rendues possibles et comparer les conditions
de leur apparition, sans exiger un accord universel sur l'objet.

La construction peut commencer avant que le sujet sache où elle mène. Un objet
retient d'abord l'attention ; une différence se détache ; une question devient
formulable ; une première réponse modifie ensuite la question. Cette suite ne
forme pas toujours une progression régulière. Le sujet peut revenir en arrière,
abandonner une piste ou comprendre que le problème était mal posé. L'enquête
devra décrire ces parcours au lieu de réduire l'intérêt à la durée de
l'attention ou au plaisir déclaré.

J'appelle \emph{intérescence} le processus temporel ainsi ouvert : rupture avec
l'indifférence, premières distinctions, constructions successives, puis
abandon, nouvelle construction ou transmission. \emph{L'intéressant} qualifie
le pouvoir de déclenchement ; \emph{l'intérescence} désigne ce qui arrive
ensuite. Le terme reste une proposition. Il ne sera conservé que s'il aide à
distinguer le déclenchement, la trajectoire et le passage d'un sujet à un autre.

\subsection{De l'objet à la réification}

Le mot \emph{objet} peut désigner tout ce qui se présente à un sujet et devient
matière à perception, à question ou à action. Le latin \emph{objectum} signifie
« ce qui est placé devant » ; le grec \emph{problēma} évoque lui aussi ce qui
est jeté ou placé devant. Les deux mots n'ont pas la même racine, mais ils
partagent une image : quelque chose se détache et appelle une réponse.

Cela conduit à une hypothèse nouvelle : l'intéressant peut être à la racine de
la réification. Ce qui devient intéressant se détache du fond. Le sujet le
remarque, y revient, lui attribue des propriétés et commence à construire avec
lui. Quelque chose devient ainsi un objet pour lui. La réification commence
lorsque cette naissance relationnelle est oubliée et que l'objet paraît exister
par lui-même. Il faut donc distinguer l'objectivation, qui rend une chose
identifiable et manipulable, de la réification, qui efface l'histoire de cette
construction \parencite{pachet2004reification}.

Cette hypothèse concerne d'abord la réification vécue. Des institutions, des
marchés, des classifications ou des programmes peuvent aussi imposer des objets
à des personnes qui ne leur portaient aucun intérêt. La programmation par
objets fournit ici une comparaison issue de la pratique de construction des
systèmes, pas une filiation historique.

\section{Propriétés à éprouver}

La définition reste une hypothèse de travail. Les propriétés suivantes serviront
à comparer les cas et à construire des objections :

\begin{enumerate}
  \item \textbf{Dépendance au sujet et au moment.} L'intéressant varie avec la
        forme, le sujet, son histoire et la situation de la rencontre.
  \item \textbf{Passage à l'activité.} Il peut conduire de l'indifférence à
        l'attention, puis à une question, une émotion ou une action qui modifie
        le sujet.
  \item \textbf{Antériorité sur la question.} Un objet peut retenir l'attention
        avant que le sujet sache ce qu'il cherche. Formuler la bonne question
        peut être la première construction.
  \item \textbf{Écart assimilable.} Il faut assez de différence pour appeler un
        travail et assez de familiarité pour que ce travail puisse commencer.
        Le déjà connu tend vers la trivialité ; l'étranger absolu vers le bruit.
  \item \textbf{Production.} Une forme intéressante suscite des comparaisons,
        hypothèses, questions ou variations qu'elle n'énonce pas elle-même.
  \item \textbf{Incomplétude structurée.} Une lacune intéresse lorsqu'elle donne
        aussi des indices pour construire ce qui manque ou reformuler le
        problème.
  \item \textbf{Transformation.} Une construction forte modifie les questions
        formulables, les distinctions disponibles ou les actions concevables.
  \item \textbf{Non-monotonie.} Plus d'information, de surprise ou d'attention
        ne produit pas nécessairement plus d'intérêt. Une explication peut
        aider la construction ou fermer trop vite l'indétermination.
  \item \textbf{Épuisement ou relance.} Certaines énigmes perdent leur intérêt
        une fois résolues ; d'autres font apparaître de nouvelles questions.
  \item \textbf{Neutralité affective.} L'intéressant peut être agréable,
        pénible, inquiétant ou répugnant. Il ne vaut pas approbation.
  \item \textbf{Transmission indirecte.} On ne transmet pas l'intérêt comme une
        propriété d'une chose. On peut préparer son apparition en installant
        une attente, en faisant voir une différence ou en fournissant un point
        de départ.
  \item \textbf{Auto-application.} Le concept de l'intéressant peut lui-même
        provoquer les distinctions et les révisions qu'il décrit.
\end{enumerate}

La difficulté ressentie devra recevoir une attention particulière. Une tâche
trop facile est déjà résolue ; une résistance sans point de départ décourage la
construction. La \emph{difficultuosité}, introduite dans \emph{La virtuosité à
la portée des caniches}, désigne la relation entre un sujet, ses compétences et
une performance perçue \parencite{pachet2027virtuosite}. La
\emph{chalépodoxie} désigne un autre phénomène : la réputation sociale d'une
difficulté peut servir à impressionner sans qu'aucune maîtrise soit acquise.
Durée d'attention, opacité et prestige ne suffisent donc pas à prouver qu'une
construction a lieu.

Ces propriétés décrivent ensemble un « épluchage » de l'objet. Une première
distinction dissipe une opacité, mais révèle souvent une nouvelle couche ou une
nouvelle résistance. La distinction acquise devient alors un moyen pour
l'opération suivante. Le processus peut s'arrêter parce que l'objet est épuisé,
parce que le sujet manque de temps ou de moyens, parce que la résistance ne
laisse plus aucun point de départ, ou simplement parce qu'il décide de finir.
Une construction achevée n'épuise donc pas toujours l'intérêt de l'objet.

\subsection{Peut-on rendre tout intéressant ?}

Rendre une matière intéressante, ce n'est pas déposer dans l'esprit d'autrui un
intérêt déjà formé. C'est accompagner le travail par lequel cette personne
commence à voir des problèmes, des différences et des questions là où elle ne
voyait qu'une matière indifférente. Cet accompagnement demande de l'ingéniosité
au sens de l'ingénieur : le professeur construit les conditions d'une
construction que l'élève devra accomplir lui-même.

Gilles de La Ménardière racontait ainsi comment son professeur avait fini, contre
toute attente, par lui rendre le droit administratif intéressant. Devenu
conseiller d'État, il faisait à son tour partager son enthousiasme pour le
problème du respect de la loi. Ce cas ne prouve pas que tout peut devenir
intéressant. Il suggère un critère : l'enseignement réussit lorsque l'intérêt
survit à la présence du professeur, que l'élève poursuit seul les questions
découvertes et qu'il peut ensuite les faire découvrir.

\section{Temps, apprentissage et récursivité}

L'intérêt change avec l'apprentissage. Une régularité opaque peut devenir
accessible, produire un gain de compréhension, puis perdre son pouvoir lorsque
le sujet la maîtrise. Les théories du progrès de compression et de la motivation
intrinsèque décrivent une version de cette évolution
\parencite{schmidhuber2009compression,oudeyer2007intrinsic}. Les situations déjà
maîtrisées et les situations inapprenables ont en commun de ne plus offrir de
progrès ; entre les deux se trouvent des zones où une amélioration est possible.
Cette idée rejoint la zone intermédiaire entre ennui et anxiété, sans réduire
l'intéressant au seul état de \emph{flow}.

Cette zone dépend du parcours. Une œuvre rencontrée trop tôt peut rester opaque,
puis devenir intéressante après l'acquisition d'une distinction ou d'une
technique. Elle peut ensuite sembler évidente et ne plus provoquer le même
travail. L'objet n'a pas changé, mais le sujet n'est plus le même. L'expérience
répétée peut aussi produire l'effet inverse : un détail négligé lors des
premières rencontres devient le point de départ d'une nouvelle exploration.

Le langage déplace cette frontière. Un mot, un concept ou une notation rassemble
des distinctions et parfois des opérations en une unité réutilisable. Ce qui
exigeait un travail devient le point de départ d'une construction nouvelle.
Nommer ne suffit pourtant pas : un terme n'aide que s'il peut être déplié en
comparaisons, prédictions, opérations ou exemples contrôlables.

La construction transforme aussi le sujet. La relation possède donc une boucle
temporelle :
\[
 I(F,S\mid H,t)\longrightarrow \text{construction}\longrightarrow S'
 \longrightarrow I(F,S'\mid H,t+1).
\]
Le même objet est ensuite rencontré par un sujet dont les attentes et les
capacités ont changé. La formule « l'intéressant est intéressant » décrit un
cas particulier de cette boucle : comprendre pourquoi une chose intéresse peut
devenir à son tour un objet de recherche. Cette fécondité ne garantit toutefois
pas la vérité de la théorie.

\section{Comment expliquer cette relation ?}

L'analyse devra répondre à trois questions concrètes. Comment l'histoire d'un
sujet change-t-elle ce qu'il remarque et ce qu'il peut faire ? Comment la
mémoire et l'attente organisent-elles le temps de l'expérience ? Comment une
forme d'ensemble produit-elle un effet que ses éléments isolés n'expliquent pas ?
La philosophie du sujet, la phénoménologie du temps et l'étude des systèmes
complexes fourniront des réponses à comparer
\parencite{petitot1993phenomenologie,petitot2002naturaliser}.

Spinoza fournit un repère important. Il refuse de traiter les affects humains
comme un domaine séparé de la nature \parencite{spinoza1966ethique}. Les idées
des affections enveloppent à la fois le corps du sujet et le corps extérieur qui
l'affecte. Ce point aide à expliquer pourquoi l'intéressant ne se trouve ni dans
la forme seule ni dans un état intérieur isolé. L'effet dépend de la rencontre,
de l'histoire du sujet et des circonstances.

L'explication comparera donc ce que les sujets racontent, ce qu'ils font, ce
qu'ils apprennent, les propriétés des formes rencontrées et les modèles qui
tentent de les décrire. Le plaisir de découvrir peut accompagner la
construction, mais il ne la définit pas : une compréhension pénible peut rester
intéressante, tandis qu'une explication agréable peut donner une impression de
clarté sans produire de distinction nouvelle.

\section{Domaines étudiés et questions communes}

\subsection{Musique}

La musique rend visibles les rapports entre mémoire, attente, surprise,
résolution, répétition et épuisement. Meyer et Narmour ont montré que l'affect
musical dépend d'attentes apprises, réalisées, retardées ou déjouées
\parencite{meyer1956emotion,narmour1990basic}. L'enquête décrira des trajectoires
d'attention à plusieurs échelles : détail mélodique ou harmonique, forme globale,
histoire des écoutes et activité de l'auditeur.

Le terme « Jotney », que j'ai créé, désigne ici des chansons qui maintiennent
une mélodie autonome et chantable tout en faisant voyager l'harmonie sans
exhiber le procédé. Elles fourniront un cas directeur. \emph{Histoire d'une
oreille} apporte deux moyens de les étudier \parencite{pachet2018oreille}.
L'écoute superposée compare mentalement des variantes. L'\emph{optale} désigne
le sentiment qu'une forme est un optimum local parce que ses variantes proches
échouent ; la \emph{pseudoptale} rappelle que la familiarité peut simuler cette
nécessité. La parcimonie tonale donnera un autre cas : une analyse compacte
devient utile si elle permet d'anticiper, de varier ou d'intervenir, mais perd sa
valeur lorsqu'elle efface une différence musicalement pertinente
\parencite{pachet2026tonalparsimony}.

La méthode consistera à comparer des passages et leurs variantes proches. Une
modification locale peut préserver la correction harmonique tout en supprimant
la tension qui appelait la suite ; une autre peut augmenter la surprise et
rendre le passage arbitraire. Ces comparaisons permettront de décrire ce que
l'auditeur anticipe, ce qu'il doit reconstruire et ce qu'il peut ensuite
réutiliser dans une écoute ou une pratique musicale. Elles aideront aussi à
distinguer l'effet d'une forme de celui de sa familiarité.

\subsection{Intelligence artificielle}

Les systèmes génératifs rendent explicites l'échantillonnage, la contrainte,
l'évaluation, la sélection et la construction de modèles. Ils montrent aussi
combien une mesure peut remplacer la question initiale. Produire une sortie
correcte, cohérente, nouvelle ou conforme à des contraintes ne suffit pas à
produire une sortie intéressante. Les métriques de surprise, de diversité, de
préférence ou de popularité restent des indicateurs locaux.

Le lien avec la musique est direct. Le Continuator et les machines de flow sont
des systèmes d'IA conçus pour interagir avec des musiciens. Ils permettent
d'étudier la relation entre un sujet et une image imparfaite, continuellement
révisée, de sa propre pratique
\parencite{pachet2008machines,addessi2005experiments}. Les meilleures imitations
immédiates ne donnent pas nécessairement les trajectoires les plus intéressantes.
Une erreur peut être rejetée ou devenir un matériau selon ce qu'elle permet de
faire ensuite. La qualité de l'interaction compte alors davantage que celle
d'un résultat isolé.

L'évolution de ces systèmes montre un déplacement. Les premiers modèles isolaient
surtout des régularités musicales ; l'écoute active a introduit le
parcours de l'utilisateur ; les machines de flow ont organisé une boucle entre
le sujet et une image calculée de son jeu. Chaque étape a rendu certaines
opérations plus visibles, mais aussi montré la limite de la représentation
précédente. La thèse ne transformera pas cette histoire en doctrine unitaire.
Elle suivra les déplacements de la question au cours de cette histoire.

Les travaux sur la \emph{Hit Song Science} donnent un résultat négatif utile :
sur 32 978 titres, des descripteurs audio et des étiquettes humaines ne classent
pas la popularité au-delà du hasard \parencite{pachetroy2008hitsong}. L'exposition
et l'influence sociale participent en effet à la production du succès
\parencite{pachet2011hitsong}. La popularité ne peut donc pas servir de mesure
simple de l'intéressant.

La thèse traitera chaque formalisation comme une hypothèse locale. Elle séparera
la cible, l'indicateur choisi et le comportement effectif du système. Les cas où
le calcul réussit tout en effaçant une différence pertinente seront conservés :
ils indiquent où la représentation doit être corrigée.

Le résultat recherché ne sera donc pas un générateur universel d'objets
intéressants. Un modèle plus modeste pourra déjà comparer les effets d'une même
forme selon l'histoire des sujets, prévoir qu'un intérêt va s'épuiser après un
apprentissage, ou expliquer pourquoi une production techniquement réussie ne
permet aucune construction nouvelle. Les échecs de prédiction auront ici autant
d'importance que les réussites s'ils obligent à préciser la relation étudiée.

\subsection{Littérature}

La littérature montre comment le langage, la narration et le style orientent
l'interprétation sans l'épuiser. Proust permettra de suivre des intérêts qui se
forment et se transforment ; \emph{Bartleby}, leur déplacement du personnage
vers le narrateur et le lecteur ; le \emph{Livre de Job}, une motivation morale
intense qui ne devient pas nécessairement intérêt exploratoire chez le
personnage \parencite{jps1985tanakh}. Ces comparaisons empêcheront de rebaptiser
toute motivation « intérescence ».

La littérature permettra aussi de distinguer une opacité structurée d'une
promesse vide de profondeur. Un texte peut multiplier les commentaires parce
qu'il fournit des indices et résiste aux paraphrases simples ; il peut aussi
entretenir l'attente d'une clef qui n'arrive jamais. La comparaison de
paraphrases, de continuations et de lectures concurrentes montrera si
l'interprétation produit des distinctions transférables ou seulement une
fascination.

\subsection{La créativité dans ces trois domaines}

La créativité distingue production du nouveau et production de l'intéressant.
Une forme peut être inédite mais arbitraire ; une autre peut employer des
matériaux connus et changer l'espace des possibles. Une solution intéressante
fait parfois découvrir après coup le problème auquel elle répond. Cette
reconstruction devra permettre des comparaisons ou des variations : elle ne
peut pas être une justification inventée après chaque réussite
\parencite{pachet2026impossibilite}.

Cette appropriation active devra être séparée de la simple satisfaction d'avoir
fabriqué quelque chose. Participer à la construction d'un objet peut augmenter
sa valeur aux yeux du sujet, comme le montre l'effet IKEA
\parencite{norton2012ikea}. Ce résultat ne prouve ni l'apprentissage ni la
justesse de l'interprétation. Il invite seulement à comparer la valeur accordée
au produit avec les distinctions réellement acquises pendant sa construction.

\subsection{La philosophie comme approche générale}

La philosophie fournit l'approche générale de la thèse. L'analyse conceptuelle
précisera l'intéressant, ses présupposés et ses contre-exemples. Elle permettra
aussi d'étudier l'auto-application : une théorie de l'intéressant doit rendre
compte des questions qu'elle fait naître, tout en distinguant cette fécondité de
sa vérité. Le projet \emph{Snarky}, qui formalise la troisième partie de
l'\emph{Éthique} de Spinoza, servira de cas. La formalisation n'assure pas une
interprétation juste, mais oblige à reconstruire les dépendances entre
définitions, axiomes et propositions.

\section{Méthode, corpus et résultats attendus}

La recherche partira de rencontres effectivement éprouvées : ce qui retient
l'attention, la résistance rencontrée, ce que le sujet commence à faire, puis
l'abandon ou la poursuite. Elle comparera ensuite les cas, formulera des
distinctions et cherchera des contre-exemples. Le concept devra enfin revenir
aux cas qui l'ont motivé afin de montrer ce qu'il explique et ce qui lui
résiste.

Les récits à la première personne auront une place précise. Ils permettront de
repérer le moment où une différence apparaît, les essais accomplis et les
changements de question. Ils ne vaudront pas à eux seuls comme preuve générale.
Ils seront comparés aux propriétés des formes, aux comportements observables,
aux variations construites et, lorsqu'ils existent, aux résultats
expérimentaux. Aucun de ces matériaux ne suffira isolément.

Le travail réunira :

\begin{itemize}
  \item l'analyse conceptuelle de l'intéressant et des notions voisines ;
  \item la lecture directe de Garve et la reconstruction de sa réception ;
  \item l'étude de cas musicaux, littéraires, scientifiques et philosophiques ;
  \item la comparaison de variantes, de continuations, de paraphrases et
        d'imitations ratées ;
  \item l'analyse de systèmes construits et, lorsque cela est pertinent,
        d'expériences sur l'attention, l'apprentissage et l'évolution des
        jugements.
\end{itemize}

Trois niveaux resteront distincts : les résultats scientifiques et les
propriétés effectives des systèmes ; les propositions philosophiques ; les
comparaisons qui montrent ce qu'un cas permet de comprendre sans prétendre le
démontrer. Une hypothèse gagnera en force si elle décrit une trajectoire,
identifie les conditions d'une variation de l'intérêt, sépare des phénomènes
jusque-là confondus ou permet de construire un contre-exemple.

Le corpus principal réunira l'essai de Garve, \emph{Histoire d'une oreille}, les
chansons Jotney, les essais sur l'impossibilité de créer et la virtuosité, des
travaux sur les systèmes génératifs et les sources philosophiques et
scientifiques directement mobilisées. Chaque source sera reliée à la
proposition qu'elle soutient et à l'état réel de sa lecture.

Le corpus sera organisé en trois ensembles : les textes directement analysés,
les cas employés pour comparer ou objecter, et les lectures encore à mener.
Cette distinction empêchera qu'une référence seulement repérée soit présentée
comme une autorité déjà acquise. Elle permettra aussi de réduire le corpus si
un texte n'éprouve aucune propriété précise du concept.

La thèse vise cinq résultats :

\begin{enumerate}
  \item une histoire qui restitue Garve comme antécédent direct et explique
        pourquoi il n'a pas fait école ;
  \item une définition relationnelle et constructive de l'intéressant, ainsi
        qu'une description du processus d'intérescence ;
  \item un ensemble de propriétés permettant de comparer apparition,
        poursuite, transmission et épuisement de l'intérêt ;
  \item une hypothèse sur le rôle de l'intéressant dans la constitution d'un
        objet puis, parfois, dans sa réification ;
  \item une série d'épreuves dans la musique, l'IA et la littérature, avec une
        analyse de la créativité dans ces trois domaines et un travail
        philosophique sur les résultats.
\end{enumerate}

\section{Plan de la thèse}

\begin{enumerate}
  \item \textbf{Constituer le problème} : notions voisines ; Garve ; raisons
        pour lesquelles il n'a pas fait école ; histoire discontinue de
        l'intéressant ; risque de remplacer la question par un indicateur.
  \item \textbf{Construire le concept} : relation entre forme, sujet, horizon
        et moment ; intérescence ; objet et réification ; propriétés ;
        apprentissage ; récursivité.
  \item \textbf{Le concept au travail} : musique, intelligence artificielle et
        littérature ; créativité dans ces trois domaines ; analyse
        philosophique ; comparaison des résultats et révision du concept.
\end{enumerate}

\section{Travaux personnels directement liés}

\begin{itemize}
  \item \emph{Histoire d'une oreille} décrit la formation d'une écoute musicale
        et fournit plusieurs concepts employés ici \parencite{pachet2018oreille}.
  \item Les travaux sur le Continuator et les systèmes réflexifs déplacent
        l'évaluation du produit isolé vers l'interaction et l'appropriation
        créative \parencite{pachet2004flowmachines,pachet2008machines}.
  \item Les recherches sur la génération sous contraintes montrent la puissance
        des modèles formels et les différences qu'ils risquent d'effacer
        \parencite{pachet2026markov,pachet2026tonalparsimony}.
  \item \emph{De l'impossibilité de créer} et \emph{La virtuosité à la portée
        des caniches} étudient la nouveauté, la fin du processus créatif, la
        difficultuosité et la chalépodoxie
        \parencite{pachet2026impossibilite,pachet2027virtuosite}.
  \item Les travaux sur la créativité musicale, la popularité et la composition
        par feedback fournissent des résultats empiriques sur la relation entre
        forme, expérience, jugement et contexte
        \parencite{pachet2006creativity,pachetroy2008hitsong,martin2016creativeprocess}.
\end{itemize}

\printbibliography[keyword={francois-pachet},title={Travaux de François Pachet}]
\printbibliography[notkeyword={francois-pachet},title={Autres références}]

\end{document}
